%! TEX root = gambar proyeksi.tex
% the main TeX file which is intended to compile, :VimtexReload after adjustment
   % this file should be included in the main file
% :h vimtex-tex-root
\documentclass[../menggambar teknik.tex]{subfiles}
%ensure the class of docs

\begin{document}

\section{Gambar Proyeksi}%
\label{sec:Gambar Proyeksi}


\subsection{Tujuan}%
\label{sub:Tujuan}

Setelah mempelajari bahan dalam bab ini, seharusnya Anda dapat:

\begin{enumerate}
	\item Membaca dan membuat gambar proyeksi sistem Amerika.
	\item Membaca dan membuat gambar proyeksi sistem Eropah.
	\item Membedakan gambar proyeksi sistem Amerika dan Eropah.
\end{enumerate}

\subsection{Pandangan Dalam Gambar Teknik Mesin}
Pada industri permesinan, gambar yang dibuat untuk diserahkan pada pekerja/teknisi pelaksana di bengkel, haruslah dibuat dalam keadaan yang memudahkan untuk dibaca dan diinterpretasikan. Agar dapat dibaca oleh orang lain, maka gambar harus dibuat dengan memberikan pandangan yang cukup. Pandangan yang cukup disini artinya tidak kurang dan juga tidak berlebihan. Pandangan gambar yang kurang akan menyebabkan kesulitan dalam menginterpretasikan maksud gambar, demikian pula gambar yang berlebihan dalam pandangan akan menyebabkan gambar menjadi rumit, sehingga kesannya semrawut dan gambarnya menjadi tumpang tindih (over lap). Untuk itu jumlah pandangan harus dibatasi seperlunya, tetapi harus dapat memberi kesimpulan bentuk benda secara lengkap.

Dalam menyajikan pandangan gambar sebuah benda, pandangan depan adalah merupakan yang pokok, sedangkan pandangan yang lain berfungsi hanya untuk memperjelas. Dengan demikian andaikata dimungkinkan cukup pandangan depan saja, maka tidak perlu dibuat pandangan yang lain, asal gambar telah memberikan pandangan yang lengkap, yang dapat memberikan satu kesimpulan mengenai bentuk dan ukuran-ukuran bagian alat yang akan dibuat.

Agar dapat membuat pandangan gambar yang baik yaitu pandangan yang tidak berlebihan atau kurang, maka berikut ini diberikan beberapa ketentuan umum untuk memilih pandangan.
\begin{enumerate}[label=(\alph*)]
	\item Jangan menggambar pandangan lebih dari yang diperlukan untuk melukis benda.
	\item Pilihlah pandangan yang sekiranya dapat memperlihatkan bentuk benda yang paling baik.
	\item Utamakanlah pandangan dengan garis yang tidak kelihatan yang paling sedikit.
	\item Pandangan sebelah kanan lebih utama dari pandangan sebelah kiri, kecuali kalau pandangan kiri memberi keterangan yang lebih banyak.
	\item Pandangan atas lebih utama dari pandangan bawah, kecuali kalau pandangan bawah memberi keterangan yang lebih banyak.
	\item Pilihlah pandangan yang sekiranya dapat mengisi ruang gambar sebaik-baiknya.
\end{enumerate}

Pandangan dalam gambar teknik mesin kebanyakan divisualisasikan dengan menggunakan proyeksi lurus. Ada dua cara untuk menggambar proyeksi lurus, yaitu proyeksi sistem Amerika (Third Angle Projection) dan proyeksi sistem Eropah (First Angle Projection). Secara lengkap kedua proyeksi ini mempunyai enam pandangan: pandangan depan, pandangan atas, pandangan samping kanan, pandangan samping kiri, pandangan bawah dan pandangan belakang.

Seperti telah dijelaskan di atas dalam penyajiannya tidak semua pandangan ini ditampilkan. Beberapa pandangan saja mungkin sudah mencukupi, seandainya obyek yang digambar tidak komplek bisa menggunakan tiga pandangan. Untuk menyajikan gambar yang sederhana, satu atau dua pandangan gambar acapkali sudah memadai.

\subsection{Gambar Proyeksi Sistem Amerika}
Pada proyeksi sistem Amerika (Third Angle Projection = Proyeksi Sudut Ketiga), bidang proyeksi terletak diantara benda dengan penglihat yang berada di luar. Untuk memproyeksikan benda pada bidang proyeksi, seolah-olah benda ditarik ke bidang proyeksi. Dengan demikian kalau bidang-bidang proyeksi dibuka, maka pandangan depan akan terletak di depan, pandangan atas terletak di atas, pandangan samping kanan terletak di samping kanan, pandangan samping kiri terletak di samping kiri, pandangan bawah terletak di bawah, dan pandangan belakang terletak di sebelah kanan samping kanan (lihat Gambar 19).


\begin{figure}
	\begin{center}
		\includegraphics[width=.95\textwidth]{proyeksiamerika.jpg}
	\end{center}
	\caption{Proyeksi Sistem Amerika}
	\label{fig:proyeksiamerika}
\end{figure}


\subsection{Gambar Proyeksi Sistem Eropa}
Pada proyeksi sistem Eropa (Fist Angle Projection = Proyeksi Sudut Pertama), benda terletak di dalam kubus diantara bidang proyeksi dan penglihat. Untuk memproyeksikan benda seolah-olah benda tersebut di dorong menuju bidang proyeksi. Dengan demikian jika bidang proyeksi di buka, maka pandangan depan tetap, pandangan samping kanan terletak di sebelah kiri, pandangan samping kiri terletak di sebelah kanan, pandangan atas terletak di sebelah bawah, pandangan bawah terletak di atas, dan pandangan belakang terletak di sebelah kanan pandangan samping kiri (lihat Gambar 20).

\begin{figure}
	\begin{center}
		\includegraphics[width=.95\textwidth]{proyeksieropa.jpg}
	\end{center}
	\caption{Proyeksi Sistem Eropa}
	\label{fig:proyeksieropa}
\end{figure}
Dari kedua proyeksi yang telah dijelaskan di atas, nampak bahwa proyeksi sistem Amerika (Third Angle Projection = Proyeksi Sudut Ketiga) penggunaannya lebih rasional dan mudah dipahami. Atas dasar itulah proyeksi sistem Amerika pemakaiannya lebih luas dibandingkan dengan sistem Eropah. Negara-negara pantai laut Pacifik, seperti USA dan Canada, juga Jepang, Korea Selatan, Australia, dan juga Indonesia menggunakan proyeksi sistem Amerika. Untuk menunjukkan penggunaan dari kedua proyeksi tersebut dapat dilihat dari lambang proyeksi seperti terlihat pada Gambar 21 di bawah ini.

\begin{figure}
	\begin{center}
		\includegraphics[width=.95\textwidth]{lambang proyeksi.jpg}
	\end{center}
	\caption{Lambang Penunjukkan Proyeksi}
	\label{fig:lambangproyeksi}
\end{figure}




\bibliographystyle{apalike}
\bibliography{../filename.bib} % required for citecompletion, not work in mainfile
\end{document}

